\clearpage
\SourceLeaf{73}%
\pdfbookmark[1]{Göttingische gelehrte Anzeigen. 5 December 1808.}{review1808}
\begin{center}\small Göttingische gelehrte Anzeigen. 5 December 1808.\end{center}
\medskip
\enlargethispage{-9\baselineskip}
The recently published programme in which Professor \textsc{Gauss} announces his lectures for the current winter term is chiefly intended to draw astronomers’ attention to the advantages of using \textit{stellar altitudes} to determine time and latitude. These observations are now unduly neglected. There is no denying that this kind of observation, especially with reflecting instruments, is harder and calls for more practice than observations of the Sun. Yet it is equally certain that these difficulties can be overcome and that in very many cases the effort is repaid many times over. Stellar altitudes are particularly worth recommending to travelling observers and seafarers. They can often take the place of solar-altitude observations when these are thwarted, as often happens, and they offer an important advantage that is lost if one confines oneself to solar altitudes. A single solar altitude gives the time only when the latitude is already known. Conversely, the time must be assumed known if one is to determine the latitude from a solar altitude outside the meridian. Combining two solar altitudes does allow time and latitude to be determined together by known methods. But if the result is to have some precision, those observations
\SourceLeaf{74}%
must be fairly far apart in time. A good clock is therefore essential, and at sea the motion of the ship must also be taken into account. By contrast, combining the observed altitude of one star with that of a second star, which may be observed immediately afterward if desired, always allows time and latitude to be determined together with great precision, provided the stars themselves are suitably chosen. The rate of the clock need be known only with enough certainty to convert the interval between the two observations into sidereal time. The problem can at once be reduced without difficulty to solving three spherical triangles, though the numerical calculation then becomes somewhat lengthy. Professor \textsc{Gauss}, however, has developed a noticeably shorter solution by purely analytical means. This may also be seen as confirmation that everything inferred from geometrical considerations can always be found at least as simply and elegantly by analysis when it is suitably applied. For practical use as well, the present solution will be found far more convenient than the one Mr. \textsc{Kraft} gave in volume 13 of the \textit{Nova Acta Petropol.} for the same problem, subject only to the immaterial restriction that the two altitude measurements must be simultaneous. The author promises to give an even more convenient indirect solution on another occasion. The present paper also explains what to consider when choosing the stars so that the results are as accurate as possible, and illustrates the method by actually applying it to a pair of stellar altitudes measured at the local observatory. The pole elevation derived from them agrees with that established by \textsc{Tob. Mayer} to within $7''$, and the time derived from them differs by only half a second of time from that determined on the same day by another method.
\par\bigskip\centerline{\rule{.24\textwidth}{.3pt}}
